\input{../tex-inputs/latex-leseansicht-vorspann}

\section[Arthur Schnitzler an Gustav Schwarzkopf, 12. 8. 1922]{L04367 Arthur Schnitzler an Gustav Schwarzkopf, 12. 8. 1922}
\nopagebreak\mylabel{L04367v}
\rehead{ }\normalsize\beginnumbering\briefempfaengerindex{Schwarzkopf, Gustav@\textsc{Schwarzkopf, Gustav}!zzzSchnitzler, Arthur@\emph{von Arthur Schnitzler}!1922-08-121@{12. 8. 1922}|(be}
\toendnotes[C]{\smallbreak\pagebreak[2]}
\correspDesc{Versand  durch Arthur Schnitzler am 12. 8. 1922 in Berchtesgaden
\newline{}Erhalt  durch Gustav Schwarzkopf im Zeitraum [13. 8. 1922 – 17. 8. 1922?] in Wien}\toendnotes[C]{\smallbreak}
\Standort{DLA, A:Schnitzler, HS.1985.1.1897.}
\physDesc{Bildpostkarte, 432 Zeichen
\newline{}Handschrift: Bleistift, lateinische Kurrent}\toendnotes[C]{\smallbreak}\pstart{}{\pb}Hrn Guſtav Schwarzkopf\pend{}\pstart{}Wien I\oindex{I., Innere Stadt@\textbf{I., Innere Stadt}, \emph{Verwaltungsgebiet}|pw}\pend{}\pstart{}Tiefer Graben 17\oindex{Wien@\textbf{Wien}!I., Innere Stadt@\textbf{I., Innere Stadt}!Tiefer Graben 17@\textbf{Tiefer Graben 17}, \emph{Wohngebäude}|pw}\pend{}{\bigskip}
\pstart
           \noindent{}\centering{}{\pb}\textcolor{gray}{\textbf{Berchtesgaden\oindex{Berchtesgaden@\textbf{Berchtesgaden}|pw} vom Malerhügel\oindex{Malerhügel@\textbf{Malerhügel}, \emph{Erhebung}|pw}}}\pend
           \vspace{1em}
\pstart
           \raggedleft{}{\pb}12. 8. 22\pend
           \vspace{0.5em}
\pstart
           Herzliche Grüße!\pend
           
\pstart
           Seit \label{K_L04367-1v}\edtext{Donnerstg hier}{\lemma{\textnormal{\emph{Donnerstg hier}}}\Cendnote{\textnormal{Vgl. A. S.: \emph{Wiener Schnitzler}, 10. 8. 1922.}}}\label{K_L04367-1} (Hotel
                  Bellevue\oindex{Hotel Bellevue [Berchtesgaden]@\textbf{Hotel Bellevue [Berchtesgaden]}, \emph{Hotel}|pw}) nach Feldafing\oindex{Feldafing@\textbf{Feldafing}, \emph{Hauptstadt}|pw} u München\oindex{München@\textbf{München}|pw} – möchte in in dieser sehr schönen
               Landschaft \label{K_L04367-4v}\edtext{gegen 14 Tage bleiben}{\lemma{\textnormal{\emph{gegen 14 Tage bleiben}}}\Cendnote{\textnormal{Er blieb bis zum 15. 9. 1922.}}}\label{K_L04367-4} Mein \label{K_L04367-2v}\edtext{Bruder\pwindex{Schnitzler, Julius 13.\,7.\,1865 Wien – 29.\,6.\,1939 ebd.@\textsc{Schnitzler, Julius} (13.\,7.\,1865 Wien – 29.\,6.\,1939 ebd.), \emph{Chirurg}|pwv} reist mit den
                  {\pb}Seinen\pwindex{Donath, Anna 23.\,1.\,1900 Wien – 27.\,12.\,1995 Cincinnati@\textsc{Donath, Anna} (23.\,1.\,1900 Wien – 27.\,12.\,1995 Cincinnati)|pwv}\pwindex{Schnitzler, Helene 16.\,7.\,1871 Budapest – September 1941 Atlantischer Ozean@\textsc{Schnitzler, Helene} (16.\,7.\,1871 Budapest – September 1941 Atlantischer Ozean)|pwv}\pwindex{Schnitzler, Karl 11.\,10.\,1896 Wien – 12.\,9.\,1981 Sydney@\textsc{Schnitzler, Karl} (11.\,10.\,1896 Wien – 12.\,9.\,1981 Sydney), \emph{Geschäftsführer}|pwv} morgen}{\lemma{\textnormal{\emph{Bruder … morgen}}}\Cendnote{\textnormal{Die Abreise verzögerte sich auf den 14. 8. 1922.}}}\label{K_L04367-2} nach Wien\oindex{Wien@\textbf{Wien}, \emph{Verwaltungsgebiet}|pw}. Was man über die dortigen Zustände hört u liest, macht
               einen erschaudern. Hier ist’s doch immer  besser! Vielleicht schreiben Sie mir
               gelegentlich eine Zeile?\pend
           
\pstart
           Grüßen Sie Ihren guten Brud\textcolor{gray}{er}\pwindex{Schwarzkopf, Max 12.\,6.\,1857 Wien – 14.\,4.\,1928 ebd.@\textsc{Schwarzkopf, Max} (12.\,6.\,1857 Wien – 14.\,4.\,1928 ebd.), \emph{Rechtsanwalt}|pwv}\pend
           \pstart Ihr \spacefill\mbox{Arthur}\pend{}\selectlanguage{ngerman}\endnumbering\briefempfaengerindex{Schwarzkopf, Gustav@\textsc{Schwarzkopf, Gustav}!zzzSchnitzler, Arthur@\emph{von Arthur Schnitzler}!1922-08-121@{12. 8. 1922}|)be}\mylabel{L04367h}
\begin{anhang}
\end{anhang}\newcommand{\dateiname}{L04367}\newcommand{\titel}{Arthur Schnitzler an Gustav Schwarzkopf, 12. 8. 1922}\newcommand{\editorInnen}{Herausgegeben von Jahnke, SelmaMüller, Martin Anton}\input{../tex-inputs/latex-leseansicht-abspann}
